\documentclass[../../../include/open-logic-section]{subfiles}

\begin{document}

\olfileid{sth}{story}{predicative}

\olsection{વિધેયક અને અવિધેયક}

રસેલનો ગણ $R$, $\Setabs{x}{x \notin
x}$ દ્વારા વ્યાખ્યાયિત થયો હતો. વધુ સ્પષ્ટ રીતે કહીએ તો $R$
બરાબર સ્વસભ્ય ન હોય એવા બધા ગણોને સમાવતો ગણ હોત. તેથી $R$ની
વ્યાખ્યા આપતી વખતે આપણે એવા ક્ષેત્ર પર પરિમાણન કરીએ છીએ જેમાં
$R$ પણ હોત (જો તે અસ્તિત્વ ધરાવતો હોત).
\sourcecorrection{OLMD-112}{મૂળમાં સ્વસભ્ય ન હોય એવા ગણોને બદલે
``સ્વસભ્ય ન હોય એવા નથી'' એવા ગણો લખાયા છે. આ વધારાનો નિષેધ
વ્યાખ્યાને ઊલટાવે છે. અહીં આપેલી $x \notin x$ વ્યાખ્યા મુજબ
સ્વસભ્ય ન હોય એવા ગણો લખ્યા છે.}

આ \emph{અવિધેયક} વ્યાખ્યા છે. વધુ સામાન્ય રીતે, કોઈ વ્યાખ્યા
અવિધેયક છે ત્યારે અને માત્ર ત્યારે જ્યારે તે વ્યાખ્યાયિત થતી
વસ્તુને સમાવતા ક્ષેત્ર પર પરિમાણન કરે, એમ કહી શકીએ.

વિરોધાભાસો પછી વ્હાઇટહેડે, રસેલે, પ્વાંકારેએ અને વેઇલે આવી
અવિધેયક વ્યાખ્યાઓને ``દુષ્ચક્રવાળી'' ગણીને નકારી હતી:
\begin{quote}
ટાળવાના વિરોધાભાસોનું વિશ્લેષણ દર્શાવે છે કે તેઓ બધા એક પ્રકારના
દુષ્ચક્રમાંથી પરિણમે છે. આવા દુષ્ચક્રો એ માન્યતામાંથી ઊભા થાય છે
કે વસ્તુઓનો સંગ્રહ એવા સભ્યો ધરાવી શકે જેમની વ્યાખ્યા માત્ર આખા
સંગ્રહ દ્વારા જ આપી શકાય[\ldots. \textparagraph]

અયોગ્ય સમગ્રતાઓ ટાળવા માટેનો સિદ્ધાંત આ રીતે કહી શકાય:
`જેમાં કોઈ સંગ્રહની \emph{બધી} વસ્તુઓ સામેલ હોય તે પોતે તે
સંગ્રહની વસ્તુ ન હોવી જોઈએ'; અથવા ઊલટી રીતે:
`જો કોઈ સંગ્રહનું સમગ્ર અસ્તિત્વ ધરાવતું હોય તો તેના કેટલાક
સભ્યોની વ્યાખ્યા માત્ર તે સમગ્ર દ્વારા જ આપી શકાય, તો તે
સંગ્રહનું કોઈ સમગ્ર અસ્તિત્વ ધરાવતું નથી.' આપણે આને
`દુષ્ચક્ર સિદ્ધાંત' કહીશું, કારણ કે તે અયોગ્ય સમગ્રતાઓની
ધારણામાં રહેલાં દુષ્ચક્રો ટાળવા દે છે.
\citep[p.~37]{WhiteheadRussell1910}
\end{quote}
અવિધેયક વ્યાખ્યાઓ નકારવામાં તેમનું અનુસરણ કરીને
\emph{દુષ્ચક્ર સિદ્ધાંત} સ્વીકારીએ, તો
(\olref[sth][story][rus]{sec}ના) વિનાશક સહજ ગણરચના-પ્રરૂપને
નીચેના જેવા પ્રરૂપથી બદલવાનો પ્રયાસ કરી શકીએ:
\sourcecorrection{OLMD-113}{મૂળમાં દુષ્ચક્ર સિદ્ધાંતને નકારવાની
વાત છે. પરંતુ પહેલાંના અવતરણમાં આ સિદ્ધાંત અવિધેયક વ્યાખ્યાઓ
નકારવાનો આધાર છે, અને પછીની વિધેયક ગણરચના તે આધારને અનુસરે છે.
તેથી અહીં સિદ્ધાંતને સ્વીકારવાની વાત કરી છે.}

\begin{defish}
\emph{વિધેયક ગણરચના.} માત્ર ગણો પર જ પરિમાણન કરતા દરેક સૂત્ર $\phi$ માટે: ગણ$^\prime$ $\Setabs{x}{\phi(x)}$ અસ્તિત્વ ધરાવે છે.
\end{defish}

જ્યાં સુધી ગણો$^{\prime}$ ગણો નથી, ત્યાં સુધી પરસ્પરવિરોધ ઊભો નહીં થાય.

દુર્ભાગ્યે વિધેયક ગણરચના બહુ \emph{વ્યાપક} નથી. છેવટે તે આપણને
નવી સત્તાઓ, ગણો$^\prime$, આપે છે. તેથી ગણો$^\prime$ પર પરિમાણન
કરતાં સૂત્રોનો વિચાર કરવો પડશે. જો તેઓ હંમેશાં ગણ$^\prime$ આપે,
તો માત્ર સ્વસભ્ય ન હોય એવા બધા ગણો$^\prime$ના ગણ$^\prime$નો
વિચાર કરવાથી રસેલનો વિરોધાભાસ ફરી ઊભો થશે. તેથી એ જ વિચાર
આગળ વધારતાં કહેવું પડે કે ગણો$^\prime$ પર પરિમાણન કરતું સૂત્ર
તેને અનુરૂપ ગણ$^{\prime\prime}$ આપે છે. પછી
ગણો$^{\prime\prime\prime}$,
ગણો$^{\prime\prime\prime\prime}$ વગેરે જોઈએ.
પ્રાઇમ-સંકેતોનો ભરાવો ટાળવા માટે તેમને ગણો$_0$, ગણો$_1$, ગણો$_2$,
ગણો$_3$, ગણો$_4$,\ldots તરીકે વિચારવું સરળ રહેશે. અને આ આપણને
પ્રકારોના (સરળ) સિદ્ધાંત તરફ દોરી જશે.

આવા સિદ્ધાંત સામે થોડા સ્પષ્ટ વાંધા છે (જોકે તેઓ
\emph{નિર્ણાયક} વાંધા છે તે સ્પષ્ટ નથી). ટૂંકમાં: મળતો સિદ્ધાંત
વાપરવામાં બોજારૂપ છે; તે જુદા પ્રકારની વસ્તુઓનું અતિશય અસ્તિત્વ
માને છે; અને છેવટે અવિધેયક વ્યાખ્યાઓ ખરેખર \emph{એટલી ખરાબ}
છે કે નહીં તે પણ સ્પષ્ટ નથી.

છેલ્લો મુદ્દો સ્પષ્ટ કરવા માટે \citeauthor{Ramsey1925}ની આ
ટિપ્પણીનો વિચાર કરો:
\begin{quote}
આપણે કોઈ પુરુષને સમૂહનો સૌથી ઊંચો પુરુષ કહી શકીએ છીએ. આમ તેને
એવી સમગ્રતા દ્વારા ઓળખીએ છીએ જેનો તે પોતે સભ્ય છે, છતાં તેમાં
કોઈ દુષ્ચક્ર નથી. \citep{Ramsey1925}
\end{quote}
રામ્સીનો મુદ્દો એ છે કે ``સમૂહનો સૌથી ઊંચો પુરુષ'' \emph{અવિધેયક}
વ્યાખ્યા છે; છતાં તે સ્પષ્ટ રીતે સંપૂર્ણપણે સ્વીકાર્ય છે.

કોઈ એવો જવાબ આપી શકે કે આ કિસ્સામાં સૌથી ઊંચી વ્યક્તિને
\emph{વિધેયક} રીતથી પણ ઓળખી શકીએ. દાખલા તરીકે, કદાચ તે પુરુષ
તરફ માત્ર આંગળી ચીંધી શકીએ. તેથી અવિધેયક વ્યાખ્યાઓ સામેનો વાંધો
સ્પષ્ટ રીતે એવી સત્તાઓ પૂરતો મર્યાદિત કરવો પડે જેમને
\emph{માત્ર} અવિધેયક રીતે જ ઓળખી શકાય. પરંતુ પછી પણ આવી
``અનિવાર્ય અવિધેયકતા'' શા માટે એટલી ખરાબ હોય, તે વિશે વધુ
સમજૂતી જોઈએ.\footnote{વધુ માટે \citet{Linnebo2010} જુઓ.}

ખરું છે કે જો વ્યાખ્યાઓ આપણને કોઈ વસ્તુની \emph{રચના કરવાની}
રીત આપે એવું ઇચ્છતા હોઈએ, તો અવિધેયક વ્યાખ્યાઓ ભારે મુશ્કેલી
ઊભી કરે છે. કારણ કે અવિધેયક વ્યાખ્યા મળતાં ખરેખર દુષ્ચક્રમાં
ફસાઈએ: અવિધેયક રીતે નિર્દિષ્ટ વસ્તુ રચવા માટે \emph{પહેલાં}
બધી વસ્તુઓ (અવિધેયક રીતે નિર્દિષ્ટ વસ્તુ સહિત) રચવી પડે, કારણ
કે તે વસ્તુ બધી વસ્તુઓના પદોમાં નિર્દિષ્ટ થાય છે. તેથી તે વસ્તુ
રચી હોય તે પહેલાં જ તેને રચવી પડે. પરંતુ ફરીથી, જો ગણોને આપણે
\emph{રચીએ} તેવી વસ્તુઓ તરીકે વિચારીએ તો જ આ ``અનિવાર્ય રીતે
અવિધેયક'' નિર્દિષ્ટ ગણો સામે ગંભીર વાંધો બને છે. અને આપણે
(કદાચ) એમ વિચારતા નથી.

આથી, સારું હોય કે ખરાબ, જે અભિગમ સામાન્ય બન્યો તે વિધેયકતા કે
અવિધેયકતા અંગે કડક વલણ લેતો નથી. તેના બદલે તે હવે સંચયી-પુનરાવર્તી
અભિગમ તરીકે ઓળખાતા અભિગમને અનુસરે છે. છેવટે આ ગણોને
``તબક્કાઓ''માં ગોઠવવા દેશે---વિધેયક અભિગમ સત્તાઓને ગણો$_0$,
ગણો$_1$, ગણો$_2$, \ldotsમાં ગોઠવે છે તેના જેવું \emph{થોડુંક}---
પરંતુ તેમની વચ્ચે પ્રકારનો કોઈ ભેદ માનીશું નહીં.

\end{document}
